\documentclass[11pt]{article}
\usepackage[a4paper,margin=1in]{geometry}
\usepackage{lmodern}
\usepackage[T1]{fontenc}
\usepackage{amsmath,amssymb}
\usepackage{xcolor}
\usepackage{booktabs}
\usepackage{array}
\usepackage{enumitem}
\usepackage{listings}
\usepackage{microtype}
\usepackage[colorlinks=true,linkcolor=blue!60!black,urlcolor=blue!60!black]{hyperref}
\usepackage{titlesec}

\definecolor{codebg}{RGB}{245,246,248}
\definecolor{codeframe}{RGB}{210,214,220}
\definecolor{accent}{RGB}{180,90,20}

\lstset{
  basicstyle=\ttfamily\small,
  backgroundcolor=\color{codebg},
  frame=single,
  rulecolor=\color{codeframe},
  framesep=6pt,
  breaklines=true,
  columns=fullflexible,
  keepspaces=true,
  xleftmargin=4pt,
  xrightmargin=4pt,
  aboveskip=8pt,
  belowskip=8pt,
}

\titleformat{\section}{\Large\bfseries\color{accent}}{\thesection}{0.6em}{}
\titleformat{\subsection}{\large\bfseries}{\thesubsection}{0.6em}{}

\setlength{\parskip}{0.5em}
\setlength{\parindent}{0pt}

\newcommand{\code}[1]{\texttt{#1}}

\title{\bfseries A Plain-Language Guide to Stochastic Differential Equations \\[2pt] and the Hull--White Model}
\author{For the p3 rates challenge: calibrating Hull--White to price a Bermudan swaption}
\date{No prior stochastic-calculus background assumed --- we build everything from scratch.}

\begin{document}
\maketitle

\section*{Part 0 --- The 30-second summary}
\addcontentsline{toc}{section}{Part 0 --- The 30-second summary}

Interest rates wiggle randomly over time. To price anything that depends on future
rates (like a swaption), we need a \textbf{mathematical model of how the rate wiggles}.
The Hull--White model is one such model. It says:

\begin{quote}
\itshape The short interest rate drifts toward a moving target, is pulled back when it
strays (mean reversion), and gets kicked around by random noise each instant.
\end{quote}

Written as an equation:
\[
  dr(t) = \underbrace{\big[\,\theta(t) - a\,r(t)\,\big]\,dt}_{\text{drift}}
        \;+\; \underbrace{\sigma(t)\,dW(t)}_{\text{randomness}}
\]

Everything below explains what every symbol means, \emph{why} it is built this way,
and how it lets you price the Bermudan in your task.

\section{What is a Stochastic Differential Equation (SDE)?}

\subsection{Start with an ordinary differential equation (ODE)}
You already know equations like $dx/dt = 5$, meaning ``$x$ grows at 5 units per
second.'' This is \textbf{deterministic}: if you know where you start, you know
\emph{exactly} where you will be at every future time, $x(t) = x(0) + 5t$. No surprises.
We often write the same thing as $dx = 5\,dt$: ``in a tiny time step $dt$, $x$ changes
by $5\,dt$.'' It is just the ODE rearranged.

\subsection{The problem: the real world is noisy}
Interest rates do not move smoothly and predictably --- they jiggle. A purely
deterministic equation can never capture a rate that is partly random. We need to
\textbf{add a random term}. That is the entire idea of an SDE:
\[
  dx = (\text{predictable part})\,dt + (\text{random part}) \times (\text{random noise})
\]
\begin{itemize}[leftmargin=1.4em]
  \item The \textbf{predictable part} is the \textbf{drift} --- the average direction
        of motion, same role as the right-hand side of an ODE.
  \item The \textbf{random part} is the \textbf{diffusion} (or volatility) --- it
        controls \emph{how big} the random jiggles are.
  \item The \textbf{random noise} itself is the new ingredient: Brownian motion.
\end{itemize}

\subsection{Brownian motion $W(t)$ --- the engine of randomness}
\textbf{Brownian motion} (a \emph{Wiener process}, hence $W$) is the standard model
of ``pure random wiggling.'' Picture a speck of dust bumped by water molecules. Its
key properties:
\begin{enumerate}[leftmargin=1.6em]
  \item \textbf{Starts at zero:} $W(0)=0$.
  \item \textbf{Random increments:} over a step $\Delta t$, the change
        $\Delta W = W(t+\Delta t)-W(t)$ is drawn from a \textbf{normal (bell-curve)}
        distribution with mean $0$ and variance $\Delta t$. In symbols
        $\Delta W \sim \mathcal{N}(0,\Delta t)$, i.e.\ $\Delta W = \sqrt{\Delta t}\,Z$
        where $Z$ is a standard normal random number.
  \item \textbf{Independent increments:} what happens in one interval is independent
        of the past. No memory.
\end{enumerate}

The crucial, slightly weird fact: because the standard deviation of a step is
$\sqrt{\Delta t}$, randomness shrinks \emph{slower} than drift as $\Delta t \to 0$.
Over a tiny step, drift contributes $\sim\!\Delta t$ (tiny) but the random part
contributes $\sim\!\sqrt{\Delta t}$ (bigger). \textbf{Over short horizons randomness
dominates; over long horizons drift dominates.} That is why a single day of rates
looks like pure noise but a decade shows a trend. In differential notation we write
$dW(t)$: an infinitesimal random kick, normal with mean $0$ and variance $dt$.

\subsection{Putting it together --- a general SDE}
\[
  dx(t) = \mu(x,t)\,dt + \sigma(x,t)\,dW(t)
\]
where $\mu(x,t)$ is the \textbf{drift}, $\sigma(x,t)$ the \textbf{diffusion/volatility},
and $dW(t)$ the Brownian kick. \textbf{How to read it:} ``Over the next instant $dt$,
$x$ moves by $\mu\,dt$ (the predictable nudge) plus $\sigma$ times a random normal kick
of size $\sqrt{dt}$.'' A simulatable version --- literally how you would code it (the
\textbf{Euler--Maruyama} scheme):

\begin{lstlisting}[language=Python]
x_next = x + mu(x, t)*dt + sigma(x, t)*sqrt(dt)*np.random.randn()
\end{lstlisting}

That one line is the computational heart of every Monte Carlo rate simulation,
including Longstaff--Schwartz in your Task~3.

\subsection{A word on It\^o calculus (just enough)}
Ordinary calculus breaks slightly when $dW$ is involved, because $dW$ is so jagged
that $(dW)^2$ is \emph{not} negligible --- it behaves like $dt$. The rulebook for
calculus with $dW$ is \textbf{It\^o calculus}, and its one rule to remember is
\[
  (dW)^2 = dt, \qquad dt\,dW = 0, \qquad (dt)^2 = 0 .
\]
You rarely need to derive things by hand here --- but this is \emph{why} Hull--White
has clean closed-form bond prices: the math works out to Gaussian (normal)
distributions, which integrate nicely.

\section{Modeling interest rates: the ``short rate''}

\subsection{What is the short rate $r(t)$?}
The \textbf{short rate} $r(t)$ is the interest rate for borrowing/lending over an
\emph{infinitesimally short} period starting at time $t$ --- the ``instantaneous''
risk-free rate at each moment. It is useful because if you know how $r(t)$ behaves over
all future times, you can price \emph{any} interest-rate product by discounting along
random paths. The price today of a \textbf{zero-coupon bond} paying \$1 at time $T$ is
the expected discount factor:
\[
  P(0,T) = \mathbb{E}\!\left[\exp\!\left(-\int_0^T r(s)\,ds\right)\right].
\]
This expectation is over all random paths the short rate could take. \textbf{A
short-rate model is just a specific SDE for $r(t)$} that makes this expectation
computable.

\subsection{What we want from a good short-rate model}
\begin{enumerate}[leftmargin=1.6em]
  \item \textbf{Mean reversion} --- rates get pulled back toward a long-run level
        rather than wandering to infinity (a plain Brownian motion would drift away
        unrealistically).
  \item \textbf{Fits today's yield curve exactly} --- the model must reproduce every
        rate the market quotes today, or every price is off from the start.
  \item \textbf{Analytical tractability} --- closed-form bond and European-option
        prices, so calibration and risk are fast (180-second budget, hundreds of
        repricings for Greeks).
\end{enumerate}
Hull--White is the simplest model delivering all three --- hence the industry
workhorse, and why the task uses it.

\section{The Hull--White model, term by term}
\[
  dr(t) = \big[\,\theta(t) - a\,r(t)\,\big]\,dt + \sigma(t)\,dW(t)
\]

\subsection{The mean-reversion term: $-a\,r(t)$}
This is the heart of the model. $a$ is the \textbf{mean-reversion speed} (a positive
constant you calibrate). Look at the sign:
\begin{itemize}[leftmargin=1.4em]
  \item If $r(t)$ is \textbf{high}, then $-a\,r(t)$ is a large negative number ---
        drift pushes the rate \textbf{down}.
  \item If $r(t)$ is \textbf{low} (or negative), $-a\,r(t)$ is positive --- drift
        pushes the rate \textbf{up}.
\end{itemize}
So the rate is always pulled toward a central level, like a spring. Bigger $a$ means a
stronger, faster pull. \textbf{Intuition:} the half-life of a shock is $\ln 2 / a$.
With $a=0.03$ a shock takes $\approx 0.69/0.03 \approx 23$ years to half-decay (slow,
persistent); with $a=0.3$ the half-life is $\approx 2.3$ years (fast reversion).

\subsection{The drift target: $\theta(t)$}
$\theta(t)$ is a \textbf{time-dependent function}, not a constant. Combined with mean
reversion, the rate reverts toward the level $\theta(t)/a$ at each instant. The clever
part that makes Hull--White special: \textbf{$\theta(t)$ is not free --- it is computed
analytically so the model reproduces today's yield curve exactly.} You do not calibrate
it. The known formula is
\[
  \theta(t) = \frac{\partial f(0,t)}{\partial t} + a\,f(0,t)
              + \frac{\sigma^2}{2a}\big(1 - e^{-2at}\big),
\]
where $f(0,t)$ is today's \textbf{instantaneous forward rate} curve (which you build in
Task~1). The takeaway: $\theta(t)$ is a plug that bends the model's average path through
every market-quoted rate --- guaranteeing property~2. In practice it gets absorbed into
the bond-price formula (Section~4), so you never compute it explicitly.

\subsection{The volatility term: $\sigma(t)\,dW(t)$}
$\sigma(t)$ is the \textbf{volatility} --- how hard the random kicks hit. In your task it
is \textbf{piecewise constant}: a different constant on each interval between exercise
dates ($\sigma_1$ on $[0,T_1]$, $\sigma_2$ on $(T_1,T_2]$, etc.). This, together with
$a$, is what you \textbf{calibrate}. Why piecewise constant? Because it gives just enough
flexibility to match the \textbf{term structure of volatility} --- the market quotes
different implied vols for different expiries, and a single constant $\sigma$ cannot fit
them all. Each $\sigma_i$ is a knob for its time bucket.
\begin{itemize}[leftmargin=1.4em]
  \item Bigger $\sigma$ $\Rightarrow$ wider spread of future rates $\Rightarrow$ more
        valuable options.
  \item $\sigma$ is what swaption prices are most sensitive to $\Rightarrow$ that is your
        \textbf{Vega}.
\end{itemize}

\subsection{Why ``one-factor'' and ``Gaussian''?}
\begin{itemize}[leftmargin=1.4em]
  \item \textbf{One-factor:} a single source of randomness --- one $dW$. The entire
        yield curve is driven by the one variable $r(t)$. A simplification (real curves
        can twist in more ways), but fast and adequate for this Bermudan.
  \item \textbf{Gaussian:} because the equation is \emph{linear} in $r$ (no $r^2$, no
        $\sqrt{r}$), the solution $r(t)$ is \textbf{normally distributed}. This is the
        magic: normal distributions integrate in closed form (closed-form bond prices),
        it is why \textbf{Jamshidian's decomposition} works, and the downside is that
        rates can go \textbf{negative} (the normal has no floor). In today's low/negative-rate
        world that is often acceptable; it matches the Normal/Bachelier vol convention
        your quotes use.
\end{itemize}

\section{From the SDE to actual prices (why HW is tractable)}

\subsection{The closed-form zero-coupon bond price}
Because $r(t)$ is Gaussian, the bond-price expectation from Section~2.1 solves exactly.
Under Hull--White, the price at time $t$ (when the short rate is $r$) of a \$1
zero-coupon bond maturing at $T$ is
\[
  P(t,T) = A(t,T)\,\exp\!\big(-B(t,T)\,r\big), \qquad
  B(t,T) = \frac{1 - e^{-a(T-t)}}{a},
\]
and $A(t,T)$ is a deterministic function built from today's discount curve and
$a,\sigma$ (it contains the $\theta$ plug, so you never compute $\theta$ directly).
\textbf{Why it matters:} this single formula is the workhorse. Every node of your
trinomial tree, or every Monte Carlo path, evaluates bond prices with it. The swap value
at exercise is a sum of these bond prices --- fast and exact, no nested simulation.

\subsection{European swaptions: Jamshidian's decomposition}
A swaption is an option on a \textbf{swap}, and a swap is a bundle of cashflows, i.e.\ a
bundle of zero-coupon bonds (a \textbf{coupon bond}). An option on a bundle is normally
hard. \textbf{Jamshidian's trick} exploits the one-factor Gaussian structure: because
\emph{every} bond price $P(t,T)$ is a \textbf{monotonic} function of the single state
variable $r$ (higher $r$ $\Rightarrow$ lower bond prices, always), there is exactly one
critical rate $r^\ast$ at which the swap is exactly at-the-money. This lets you
\textbf{decompose} the swaption into a \textbf{portfolio of options on individual
zero-coupon bonds}, each with a \textbf{closed-form} Hull--White price (a Black-76-like
formula). Sum them $\Rightarrow$ the European swaption price, analytically. This is the
method for \textbf{Task~2 (calibration)}: you price European swaptions thousands of times
searching for $(a,\sigma_1..\sigma_N)$ that match market vols --- closed-form Jamshidian
makes each evaluation microseconds instead of a simulation.

\subsection{Bermudan swaptions: why we need a tree or Monte Carlo}
A \textbf{Bermudan} swaption can be exercised at \emph{several} dates (but only once). At
each exercise date you decide
\[
  V = \max\big(\text{exercise value},\ \text{continuation value}\big),
\]
where the exercise value is the value of entering the swap now (closed-form, from the
bond formula) and the continuation value is the value of keeping the option alive (it
depends on all later decisions). This is an \textbf{optimal stopping problem} --- no
single closed form --- so you solve it \textbf{backward} numerically:
\begin{itemize}[leftmargin=1.4em]
  \item \textbf{Trinomial tree:} discretize $r(t)$ onto a recombining lattice (up / stay
        / down, with the three branches calibrated to HW's mean and variance). Roll
        backward from the last date: at each node take $\max(\text{exercise},
        \text{continuation})$, discounting expected values back one step. The Gaussian
        structure is what makes the tree \emph{recombine} (up-then-down lands on the same
        node as down-then-up), keeping it small.
  \item \textbf{Monte Carlo + Longstaff--Schwartz:} simulate many random $r(t)$ paths
        (the Euler step from Section~1.4). The trouble with MC and early exercise is you
        do not know the continuation value along a path; Longstaff--Schwartz estimates it
        by \textbf{regressing} realized future payoffs against functions of the current
        state, at each exercise date, working backward. The fitted value \emph{is} the
        continuation-value estimate.
\end{itemize}
Either way, the \textbf{early-exercise premium} (the extra value of a Bermudan over a
European) falls out of this backward induction.

\subsection{Greeks: bump-and-reprice (Task 4)}
\begin{itemize}[leftmargin=1.4em]
  \item \textbf{Vega} = sensitivity to each market vol quote: bump one vol by $+1$bp,
        \textbf{re-calibrate} $(a,\sigma)$, re-price, take the difference.
  \item \textbf{Delta} = sensitivity to each zero rate: bump one zero rate by $+1$bp,
        rebuild the curve, re-calibrate, re-price, difference.
\end{itemize}
This is why \textbf{speed matters}: every Greek is a full re-calibration plus
re-pricing. Closed-form Jamshidian (calibration) plus an efficient tree (pricing) keep
you inside the 180-second budget.

\section{Jamshidian's decomposition --- step by step}

\textbf{The problem.} A payer swaption is the right to enter a swap, and a swap is
worth a \emph{bundle} of cashflows (a coupon bond). An option on a \emph{sum} of
things is normally hard, because the sum's value depends on how all the pieces move
together.

\textbf{Jamshidian's insight} (works precisely because Hull--White is one-factor):
everything at expiry depends on a single random number --- the state $x$. And every
bond price $P(\text{expiry},T\mid x)$ is \textbf{monotonic} in $x$ (higher $x$
$\Rightarrow$ higher rates $\Rightarrow$ lower bond price). So the swap's value is also
monotonic in $x$: there is \textbf{exactly one} critical state $x^\ast$ where the swap is
precisely at-the-money (worth zero). On one side of $x^\ast$ you exercise, on the other
you do not. Because that boundary is a single point, the option on the bundle
\textbf{splits into a sum of independent options on each individual bond}, each struck at
that bond's value at $x^\ast$ --- and each of those has a closed-form (Black-style) price.

\textbf{The recipe (exactly what your \code{EuroPricer} does):}
\begin{enumerate}[leftmargin=1.6em]
  \item \textbf{Write the swap value at expiry as a function of the state $x$:}
        \[
          \text{swap}(x) = \underbrace{1}_{\text{receive floating}}
            - \underbrace{\sum_i K\,\tau_i\,P(e,T_i\mid x)}_{\text{pay fixed leg}}
            - \underbrace{P(e,T_n\mid x)}_{\text{final notional}},
        \]
        which is positive when the payer is in-the-money.
  \item \textbf{Find the critical state $x^\ast$} where $\text{swap}(x^\ast)=0$. This is a
        one-dimensional root-find (your code uses \code{brentq} over $\pm 8$ standard
        deviations, with guards for the all-in / all-out cases).
  \item \textbf{Turn each cashflow into a bond option.} At $x^\ast$, each bond price
        $P(e,T_i\mid x^\ast)$ becomes a \textbf{strike} $K_i$. The swaption is now a
        \textbf{portfolio of bond put options}, one per cashflow, each weighted by its
        cashflow size $c_i = K\tau_i\,(+1$ at the final date$)$.
  \item \textbf{Price each bond put in closed form} with a Black-76-style formula using the
        bond's Hull--White volatility $v_i = B(e,T_i)\cdot\text{std}_e$, then \textbf{sum}:
        $\;\text{price} = \sum_i c_i \cdot \text{Put}_i$.
\end{enumerate}

\textbf{Why a put, for a payer?} A payer gains when rates \emph{rise}, i.e.\ when bond
prices \emph{fall}. An option that pays off when a bond price falls below a strike is a
\textbf{put} on that bond. So a payer swaption equals a basket of bond puts. (A receiver
would be a basket of bond calls.)

\textbf{Payoff:} an exact European price in microseconds, with no simulation --- which is
what makes the thousands of repricings in calibration (Task~2) feasible inside the time
limit.

\section{The backward-induction engine (your Bermudan PDE pricer)}

A Bermudan can be exercised at several dates, and the value today depends on a whole
chain of future decisions --- so you cannot price it forward. You solve it
\textbf{backward}, from the last date to today. Your code does this with a \textbf{PDE on a
grid}, which is a smooth, deterministic cousin of a tree.

\subsection{The state variable}
Instead of tracking $r$ directly, work with $x(t) = r(t) - \alpha(t)$, where $\alpha(t)$ is
the deterministic (curve-fitting) part. $x$ is a clean mean-reverting process centred at
$0$ (an \textbf{Ornstein--Uhlenbeck} process), which keeps the grid symmetric and stable.
The short rate is recovered as $r = x + \alpha(t)$.

\subsection{The grid}
\begin{itemize}[leftmargin=1.4em]
  \item \textbf{Space:} a row of possible $x$ values from $-x_{\max}$ to $+x_{\max}$ (your
        code: $401$ points, $x_{\max}\approx 6$ standard deviations --- wide enough to hold
        essentially all the probability).
  \item \textbf{Time:} small steps (daily, $1/365$) from today out to \code{swap\_end},
        with the exercise dates inserted \emph{exactly} so decisions land on grid times.
\end{itemize}

\subsection{The backward march}
\begin{enumerate}[leftmargin=1.6em]
  \item \textbf{Start at the end.} At \code{swap\_end} the option is worth $0$ everywhere
        (terminal condition $V=0$).
  \item \textbf{Step back one $dt$ at a time.} Each step answers: ``given the value one
        step in the future, what is it worth now?'' --- the discounted expected future
        value, accounting for (a)~the random spreading of $x$ (diffusion
        $\tfrac12\sigma^2 V_{xx}$), (b)~the mean-reverting drift ($-a\,x\,V_x$), and
        (c)~discounting at the current rate ($-(x+\alpha)\,V$). Your code does this with one
        \textbf{implicit finite-difference solve} per step --- a fast tridiagonal system
        (\code{solve\_banded}). ``Implicit'' just means a stable scheme that will not blow
        up even with large steps.
  \item \textbf{At every exercise date, apply the early-exercise test.} Compute the value
        of exercising now (entering the payer swap at that node's rate):
        \[
          \text{exercise}(x) = 1 - P(t,\text{swap\_end}\mid x) - K\cdot\text{annuity}(x),
        \]
        then take, at every grid point,
        \[
          V = \max\big(\text{continuation}(x),\ \text{exercise}(x)\big).
        \]
        You exercise wherever exercising beats holding --- this single $\max$ is what
        captures the \textbf{early-exercise premium}.
  \item \textbf{Keep marching} until $t=0$.
\end{enumerate}

\subsection{Read off the answer}
At $t=0$ today's state is $x=0$, so the price is $V$ interpolated at $x=0$, multiplied by
the notional.

\subsection{Why backward, and how it relates to trees / Longstaff--Schwartz}
You must go backward because the exercise decision at each date needs the value of
\emph{continuing}, which depends on all \emph{later} dates. Starting at the end and working
back means the continuation value is already known when you reach each decision --- the
essence of backward induction (and of optimal-stopping problems generally). Same idea,
different machinery: a \textbf{tree} discretizes into up/down branches; \textbf{Longstaff--Schwartz}
uses simulated paths plus regression to estimate continuation values; your \textbf{PDE}
discretizes the governing equation on a fixed grid and solves it directly. For a clean
one-factor model like this, the PDE is smooth, stable, and very accurate.

\section{The whole pipeline in one picture}
\begin{lstlisting}
  Task 1                Task 2                    Task 3              Task 4
+---------+   +----------------------+   +------------------+   +----------+
| Zero    |   | Calibrate a, sigma:  |   | Price Bermudan:  |   | Bump &   |
| curve ->| ->| choose params so HW  | ->| backward induct  | ->| reprice  |
| DFs,    |   | European prices      |   | on tree / LSMC,  |   | -> Delta,|
| forwards|   | (Jamshidian) match   |   | max(exercise,    |   |    Vega  |
|         |   | market vols (min RMSE)|  | continuation)    |   |          |
+---------+   +----------------------+   +------------------+   +----------+
   curve         the SDE gets its            the SDE is            sensitivities
   feeds theta   a and sigma pinned          simulated/lattice'd   of the price
\end{lstlisting}

\section{Glossary of every symbol}
\renewcommand{\arraystretch}{1.3}
\begin{center}
\begin{tabular}{>{\ttfamily}l l p{6.2cm} l}
\toprule
\normalfont Symbol & Name & What it is & Free or fixed? \\
\midrule
r(t)         & short rate            & instantaneous risk-free rate at $t$           & the random variable \\
dr(t)        & change in $r$         & how $r$ moves over an instant $dt$            & --- \\
a            & mean-reversion speed  & how strongly $r$ is pulled back               & \textbf{calibrated} \\
theta(t)     & drift target          & bends model to fit today's curve              & fixed (from curve) \\
sigma(t)     & volatility            & size of random kicks, piecewise-constant      & \textbf{calibrated} \\
dW(t)        & Brownian kick         & random normal noise, variance $dt$            & the randomness \\
W(t)         & Brownian motion       & cumulative random walk                        & --- \\
P(t,T)       & bond price            & value at $t$ of \$1 paid at $T$               & closed-form \\
f(0,t)       & forward rate          & today's instantaneous forward curve           & from Task 1 \\
A,B          & HW coefficients       & building blocks of $P(t,T)$                   & from $a,\sigma$, curve \\
\bottomrule
\end{tabular}
\end{center}

\section{Common pitfalls \& sanity checks}
\begin{enumerate}[leftmargin=1.6em]
  \item \textbf{Number of $\sigma$ parameters is dynamic.} It equals the number of
        breakpoint intervals $[0,T_1,\dots,T_k,\text{swap\_end}]$. Read it from the input
        every time --- never hardcode (hidden cases differ).
  \item \textbf{Negative rates are allowed} in Hull--White (Gaussian). Do not floor them
        at zero --- the Normal/Bachelier convention expects them.
  \item \textbf{$\theta(t)$ is never calibrated.} If you are fitting $\theta$, you have
        misunderstood --- it is analytic from the curve. Only $a$ and $\sigma$ are free.
  \item \textbf{Mind the metric.} Calibration minimizes RMSE of vols in bps, but pricing
        is scored on MAPE of the dollar price --- different objectives.
  \item \textbf{Recombining tree:} verify up-then-down lands on down-then-up, or the tree
        blows up exponentially and busts the time limit.
  \item \textbf{Sanity-check the bond formula:} $P(t,t)=1$, and $P(0,T)$ must match your
        Task-1 discount factors exactly. If not, $A(t,T)$ is wrong.
  \item \textbf{Greek signs:} a \emph{payer} swaption gains value when rates rise (positive
        delta to rates); vega is positive everywhere (more vol = more option value).
\end{enumerate}

\section{If you want to go deeper}
\begin{itemize}[leftmargin=1.4em]
  \item \textbf{Brigo \& Mercurio, \emph{Interest Rate Models --- Theory and Practice}}:
        the definitive reference; Chapter~3 covers Hull--White and Jamshidian in full.
  \item \textbf{Hull, \emph{Options, Futures, and Other Derivatives}}: gentler, with the
        trinomial tree construction step by step.
  \item \textbf{Original papers:} Hull \& White (1990) for the model; Jamshidian (1989)
        for the decomposition; Longstaff \& Schwartz (2001) for the MC regression method.
\end{itemize}

\vspace{1em}
\hrule
\vspace{0.8em}
\textit{The one sentence to remember: Hull--White is just an SDE that makes the short rate
spring back toward a curve-fitting target while getting randomly kicked --- and because
it is linear (hence Gaussian), it hands you closed-form bond and European prices, which
you then stack into a backward-induction pricer for the Bermudan.}

\end{document}
